\documentclass[12pt]{article}
\usepackage[utf8]{inputenc}
\usepackage[T1]{fontenc}
\usepackage{hyperref}
\title{Ngarigu Language Resources for K-12 Curriculum Development}
\author{Compiled by Claude Sonnet 5.5 (claude-sonnet-5-5) for Andr\'as Kornai}
\date{}
\begin{document}
\maketitle

\section{Introduction}
Ngarigu (also Ngarigo; ISO 639-3 xni) is a Pama-Nyungan language of the Monaro and Snowy Mountains region of south-eastern Australia. No one has grown up speaking it. Sources disagree on status: ASJP lists it as extinct, while the Endangered Languages Project includes it because people identify with it and are working to revitalise it. Documentation consists mainly of early wordlists, so any programme would be a revival programme built on limited records. The sources found here are scholarly or reference sources, and no school-ready materials were located.

\section{Primary Educational Resources}
No primers, textbooks or school programmes specific to Ngarigu were found in the search. The Endangered Languages Project mentions a revitalisation programme and a November 2025 story about Indigenous teens connecting with language on the NSW far south coast. Whether that programme covers Ngarigu specifically was not confirmed, and the search did not give its details.

\section{Linguistic and Reference Materials}
\begin{itemize}
\item Harold Koch's chapter on documentary sources for Ngarigu, in \emph{Language, Land and Song}. It covers Lhotsky's 1834 vocabulary of 171 words and other sources: \url{https://www.elpublishing.org/docs/6/01/LLS-Chapter-11-Koch.pdf}
\item Wikipedia article on Ngarigo. It lists early wordlist compilers (Lhotsky, du V\'e, Bulmer, Robinson, Howitt, Mathews) and Luise Hercus's 1963 recordings of terms from descendants: \url{https://en.wikipedia.org/wiki/Ngarigo}
\item ASJP wordlist: \url{https://asjp.clld.org/languages/NGARIGU}
\item Glottolog bibliographic record: \url{https://glottolog.org/resource/reference/id/ozbib:850}
\end{itemize}

\section{Religious and Community Resources}
No specific community organisations were identified in the search. The Endangered Languages Project pages are a starting point for finding revitalisation contacts: \url{https://endangeredlanguages.com/lang/6712?hl=de}.

\section{Digital and Audio Resources}
Resource listing at the Endangered Languages Project: \url{https://endangeredlanguages.com/node/40144/resources}. No audio, apps or corpora were confirmed in the search results.

\section{Curriculum Development Context}
Because the language is documented mainly through short wordlists, a K-12 programme could introduce place names, vocabulary and cultural context. It should be led by Ngarigu community members and Aboriginal language organisations, whom the search did not name. Material from the chapter by Koch should be used with attention to the varieties it distinguishes (Monaro, Tumut/Walgalu, Canberra-Queanbeyan, Omeo). Access to historical records should respect community protocols.
\end{document}